\section{Introduction}
\label{sec:intro}

In laboratory environments, maintaining a good level of cleanliness is often a challenge. While in modern domestic situations, an autonomous floor-cleaning robot is regularly used nowadays, in laboratory environments this is often not yet the case. In a laboratory, what actually falls onto the floor usually consists of graspable objects, while in domestic spaces these types of robots are generally made to remove dust and hair instead. Hence, a robot for laboratory-cleaning purposes should be able to grab objects instead of dust or hair. But as one may not always want to throw away all objects found in a laboratory setting, one would also want to separate trash from other fallen objects, such as electronics. Hence, the goal for this project was to develop an autonomous laboratory-cleaning robot that can detect, separate and dispose of two different types of objects, all based on the MIRTE Master robotic platform.

This report outlines the roadmap and methodology required to make such a robot. \newline Section 1 (\href{https://matt-rbt.github.io/Lab-Cleanup-Robot-using-the-Mirte-Master-Platform/materials/}{Materials}) includes information about the hardware and software used to realize the robot. \newline Section 2 (\href{https://matt-rbt.github.io/Lab-Cleanup-Robot-using-the-Mirte-Master-Platform/theory/}{Theory}) explains in detail how the theory behind the implemented ideas works. \newline Section 3 (\href{https://matt-rbt.github.io/Lab-Cleanup-Robot-using-the-Mirte-Master-Platform/systemoverview/}{System Overview}) breaks down the implementation of all major aspects of the robot into the robot. \newline Section 4 (\href{https://matt-rbt.github.io/Lab-Cleanup-Robot-using-the-Mirte-Master-Platform/methods/}{Experimental Methods}) shows how the implementations are tested. \newline Section 5 (\href{https://matt-rbt.github.io/Lab-Cleanup-Robot-using-the-Mirte-Master-Platform/results/}{Results}) covers the results of the tests that have been performed. \newline Section 6 (\href{https://matt-rbt.github.io/Lab-Cleanup-Robot-using-the-Mirte-Master-Platform/conclusion/}{Conclusion}), at the end, revisits the purpose of the project and reflects on that using the results that have been obtained and explained in Section 5 (\href{https://matt-rbt.github.io/Lab-Cleanup-Robot-using-the-Mirte-Master-Platform/results/}{Results}).